
\section{固体理论：玻尔兹曼方程与弛豫时间近似——从分布函数到电导}\label{sec:boltzmann}

\subsection*{恒定电场下费米分布的$\bm{k}$空间平移}

\begin{modelbox}[题目]
在恒定电场 $\bm{E}$ 中，弛豫时间 $\tau$ 为与 $\bm{k}$ 无关的常量。证明一级近似下的非平衡分布
\[
f = f_0 + f_1 = f_0 + \frac{e\tau}{\hbar}\,\bm{E} \cdot \nabla_{\!\bm{k}} f_0
\]
等价于将平衡分布 $f_0(\bm{k})$ 在 $\bm{k}$ 空间中整体平移了 $\Delta\bm{k} = e\tau\bm{E}/\hbar$。
\end{modelbox}

\begin{tipbox}[考点快速分析]
\begin{itemize}
    \item \textbf{玻尔兹曼方程：}分布函数 $f(\bm{r},\bm{k},t)$ 的演化方程
    \item \textbf{弛豫时间近似：}碰撞项简化为 $-(f - f_0)/\tau$
    \item \textbf{微扰展开：}$f = f_0 + f_1 + f_2 + \cdots$ 按电场幂次展开
    \item \textbf{平移算符：}$f_0(\bm{k} + \delta\bm{k}) \approx [1 + \delta\bm{k} \cdot \nabla_{\!\bm{k}}]f_0(\bm{k})$
\end{itemize}
\end{tipbox}

\subsection{玻尔兹曼方程是什么？——从零开始的输运理论}

在固体中，电子数目极其庞大（$\sim\!10^{23}\,\text{cm}^{-3}$），不可能逐个求解薛定谔方程。我们转而在相空间 $(\bm{r},\bm{k})$ 中描述电子的统计分布 $f(\bm{r},\bm{k},t)$。

\begin{derivationbox}[从牛顿到玻尔兹曼——半经典方程从哪来？]
\begin{enumerate}
    \item \textbf{速度方程} $\dot{\bm{r}} = \frac{1}{\hbar}\nabla_{\!\bm{k}}E(\bm{k})$
    
    对 Bloch 电子，波包由能带 $E_n(\bm{k})$ 描述。量子力学中可严格证明：波包中心的群速度为
    \[
    \bm{v}_g = \frac{1}{\hbar}\nabla_{\!\bm{k}}E_n(\bm{k}).
    \]
    这来自波包 $e^{i\bm{k}\cdot\bm{r}}u_{n\bm{k}}(\bm{r})$ 的相位演化：$\partial_{\bm{k}}\omega = \partial_{\bm{k}}(E/\hbar) = \bm{v}_g$。

    \item \textbf{动量方程} $\hbar\dot{\bm{k}} = -e(\bm{E} + \bm{v} \times \bm{B})$
    
    外电磁场中，电子的哈密顿量为 $\hat{H} = \frac{1}{2m}(\bm{p} + e\bm{A})^2 - e\phi$。在能带理论框架下，将 $\bm{k}$ 理解为波包中心参数，可证明（见固体物理教材的半经典动力学章节）：
    \[
    \hbar\dot{\bm{k}} = -e\bigl(\bm{E} + \bm{v} \times \bm{B}\bigr).
    \]
    注意这里的 $\bm{k}$ 是晶体动量（含规范势），不是自由粒子的动量。

    \item \textbf{无碰撞时的分布函数演化}
    
    对分布函数 $f(\bm{r},\bm{k},t)$，如果电子只受外场驱动、没有碰撞，则 $f$ 沿经典轨道守恒：
    \[
    \frac{df}{dt} = \frac{\partial f}{\partial t} + \dot{\bm{r}} \cdot \nabla_{\!\bm{r}}f + \dot{\bm{k}} \cdot \nabla_{\!\bm{k}}f = 0.
    \]
    这被称为\textbf{无碰撞玻尔兹曼方程}（或 Vlasov 方程），是刘维定理在单粒子分布函数上的体现。刘维定理的原始表述是：相空间中代表点的密度在运动过程中保持不变（$d\rho/dt = 0$）。对单粒子分布函数，等价于 $f$ 沿轨道守恒。

    \item \textbf{碰撞作为源项}
    
    杂质、声子、电子--电子散射使 $f$ 偏离自由演化。玻尔兹曼方程把碰撞写成源项：
    \[
    \frac{\partial f}{\partial t} + \dot{\bm{r}} \cdot \nabla_{\!\bm{r}}f + \dot{\bm{k}} \cdot \nabla_{\!\bm{k}}f = \left(\frac{\partial f}{\partial t}\right)_{\!\text{coll}}.
    \]
\end{enumerate}
\end{derivationbox}

\subsection{弛豫时间近似}

碰撞项的严格计算极其复杂（需要量子散射理论）。\textbf{弛豫时间近似}（Relaxation Time Approximation, RTA）用一个极其简单但物理透明的模型代替碰撞：

\begin{derivationbox}[弛豫时间近似的物理图像]
假设系统偏离平衡后，碰撞会指数地把分布拉回到平衡态 $f_0$：
\[
\left(\frac{\partial f}{\partial t}\right)_{\!\text{coll}} = -\frac{f - f_0}{\tau}.
\]

\textbf{为什么是这样？}
\begin{itemize}
    \item $\tau$ 是平均自由时间：电子两次碰撞之间的平均时间
    \item $f > f_0$ 时，右边为负，$f$ 被拉低
    \item $f < f_0$ 时，右边为正，$f$ 被抬高
    \item 如果没有外场持续驱动，$f(t) = f_0 + (f(0) - f_0)e^{-t/\tau}$，确实指数弛豫
\end{itemize}
\end{derivationbox}

\begin{warnbox}[注意]
严格来说 $\tau = \tau(\bm{k})$ 可以依赖于能量/动量方向（如声学声子散射 $\tau \propto T^{-1}\varepsilon^{-1/2}$，电离杂质散射 $\tau \propto \varepsilon^{3/2}$）。本题假设 $\tau$ 为常数，是最简单的模型。
\end{warnbox}

\subsection{稳态方程与微扰展开}

\begin{derivationbox}[稳态方程与微扰展开]
对恒定电场、均匀系统（无温度梯度，$\nabla_{\!\bm{r}}f = 0$），稳态下 $\partial f/\partial t = 0$，玻尔兹曼方程简化为：
\[
-\frac{e}{\hbar}\bm{E} \cdot \nabla_{\!\bm{k}}f = -\frac{f - f_0}{\tau}.
\]

将 $f$ 按电场 $\bm{E}$ 的幂次展开：
\[
f = f_0 + f_1 + f_2 + \cdots
\]
其中 $f_n \propto E^n$。代入方程，两边同次幂相等：
\begin{itemize}
    \item $\mathcal{O}(E^0)$：$0 = -(f_0 - f_0)/\tau$，恒成立（$f_0$ 是平衡解）
    \item $\mathcal{O}(E^1)$：$\displaystyle -\frac{e}{\hbar}\bm{E} \cdot \nabla_{\!\bm{k}}f_0 = -\frac{f_1}{\tau}$
    \item $\mathcal{O}(E^2)$：$\displaystyle -\frac{e}{\hbar}\bm{E} \cdot \nabla_{\!\bm{k}}f_1 = -\frac{f_2}{\tau}$
\end{itemize}
\end{derivationbox}

\subsection{一级近似解与$\bm{k}$空间平移}

\begin{derivationbox}[一级近似解与$\bm{k}$空间平移]
从 $\mathcal{O}(E^1)$ 方程解得：
\[
f_1 = \frac{e\tau}{\hbar}\,\bm{E} \cdot \nabla_{\!\bm{k}}f_0.
\]

一级近似下的非平衡分布：
\[
f = f_0 + f_1 = f_0 + \frac{e\tau}{\hbar}\,\bm{E} \cdot \nabla_{\!\bm{k}}f_0 = \Bigl[1 + \frac{e\tau}{\hbar}\,\bm{E} \cdot \nabla_{\!\bm{k}}\Bigr]f_0(\bm{k}).
\]

\textbf{泰勒展开}

对多元函数 $f_0(\bm{k})$，沿方向 $\delta\bm{k}$ 的一阶泰勒展开为：
\[
f_0(\bm{k} + \delta\bm{k}) = f_0(\bm{k}) + \delta\bm{k} \cdot \nabla_{\!\bm{k}}f_0 + \mathcal{O}(|\delta\bm{k}|^2).
\]

比较 $f = [1 + \frac{e\tau}{\hbar}\bm{E} \cdot \nabla_{\!\bm{k}}]f_0$ 与上式，取
\[
\delta\bm{k} = \frac{e\tau}{\hbar}\bm{E},
\]
则在一级近似下：
\[
f(\bm{k}) \approx f_0\!\left(\bm{k} + \frac{e\tau}{\hbar}\bm{E}\right).
\]
\end{derivationbox}

\begin{insightbox}[物理图像：费米球的刚性漂移]
对自由电子气，$f_0(\bm{k})$ 是费米--狄拉克分布，在低温下近似为费米球（半径 $k_F$）。

一级近似 $f(\bm{k}) = f_0(\bm{k} + \Delta\bm{k})$ 意味着：
\begin{itemize}
    \item 整个费米球在 $\bm{k}$ 空间中刚性平移了 $\Delta\bm{k} = e\tau\bm{E}/\hbar$
    \item 电子的平均动量漂移：$\langle\hbar\bm{k}\rangle = \hbar\Delta\bm{k} = e\tau\bm{E}$
    \item 平均漂移速度：$\bm{v}_d = \langle\hbar\bm{k}\rangle/m^* = e\tau\bm{E}/m^*$
    \item 电流密度：$\displaystyle \bm{j} = -ne\bm{v}_d = \frac{ne^2\tau}{m^*}\bm{E} = \sigma\bm{E}$
\end{itemize}

这正是欧姆定律的微观推导，电导率 $\displaystyle \sigma = \frac{ne^2\tau}{m^*}$（Drude 公式）。
\end{insightbox}
